\documentclass[11pt,letterpaper]{scrartcl}

% =====================================================================
% Engine: xelatex (required for the unicode monofont below).
% =====================================================================

\usepackage{fontspec}
\setmonofont{DejaVu Sans Mono}[Scale=0.85]

\usepackage[margin=1in]{geometry}
\usepackage[colorlinks=true,
            linkcolor=blue!50!black,
            urlcolor=blue!50!black,
            citecolor=blue!50!black]{hyperref}
\usepackage{booktabs}
\usepackage{longtable}
\usepackage{enumitem}
\usepackage{array}
\usepackage{xcolor}
\usepackage{listings}

% Listings: shaded background, wrap long lines, use the mono font.
\definecolor{codebg}{HTML}{F5F5F5}
\lstset{
  basicstyle=\ttfamily\small,
  backgroundcolor=\color{codebg},
  breaklines=true,
  breakatwhitespace=false,
  columns=fullflexible,
  frame=single,
  framerule=0pt,
  framesep=4pt,
  xleftmargin=4pt,
  xrightmargin=4pt,
  showstringspaces=false,
  keepspaces=true,
}

\setlist[itemize]{topsep=2pt,itemsep=2pt,parsep=0pt}
\setlist[enumerate]{topsep=2pt,itemsep=2pt,parsep=0pt}

\title{A Phased Refactor of a Personal Shell Toolbox\\
       \large An audit and remediation plan for
       \texttt{.zshrc} and \texttt{\textasciitilde/bin}}
\author{R.\ G.\ Thomas}
\date{2 May 2026}

\begin{document}

\maketitle

\begin{abstract}
\noindent
This white paper documents a concrete, phased refactor of a working
quantitative researcher's personal shell toolbox. The toolbox under
audit comprises a 736-line \texttt{.zshrc} and a
\texttt{\textasciitilde/bin} directory holding roughly sixty
top-level entries. The audit identifies two categories of
misalignment: shell functions that mutate no shell state and should
be versioned scripts; and microscripts that exist only because their
authors pre-dated discovering the alias. A guiding principle resolves
both inversions: \emph{a helper should be a shell function only when
it must mutate the calling shell, otherwise it should be a versioned,
executable script in} \texttt{\textasciitilde/bin}. The paper applies
this rule to specific helpers, sequences the refactor into nine
independently shippable phases, lists effort estimates and acceptance
criteria, and surfaces the risks worth managing during execution.
\end{abstract}

\tableofcontents

\section{Context and Scope}

A personal shell toolbox accumulates over years of incremental work.
Each helper is added in response to a specific irritation, and the
choice between placing it in \texttt{.zshrc} as a function or in
\texttt{\textasciitilde/bin} as a script is typically made by which
file happens to be open in the editor at the moment. After several
years of accretion, the toolbox remains functional, but the boundary
between the two locations has eroded: scripts in
\texttt{\textasciitilde/bin} are sometimes aliases in disguise;
functions in \texttt{.zshrc} sometimes span hundreds of lines of
program logic that have no need for the calling shell's state.

The toolbox examined here exhibits both inversions. The objective of
this refactor is not to rewrite the toolbox from scratch (which would
discard institutional knowledge embedded in working code), but to
move each existing helper into the location its actual behaviour
demands. The plan is structured so that each phase is independently
shippable and reversible; no phase commits the author to the next.

\paragraph{Out of scope.} Replacement of the toolbox by a
configuration-management system (Ansible, Nix, chezmoi) is a
distinct, larger project and is not addressed here. Migration of
Python or R helpers, whose execution model differs from POSIX
scripts, is also out of scope. The plan operates only on shell-level
artefacts.

\section{Audit Summary}

\subsection{The startup file}

\texttt{.zshrc} comprises 736 lines distributed approximately as
follows:

\begin{itemize}
  \item Configuration that legitimately belongs in the startup file
        (PATH, history, completion, aliases, plugin loading, hooks):
        roughly 200 lines.
  \item Shell-state functions (navigation, dirstack, fzf-then-cd):
        roughly 80 lines.
  \item Functions that operate exclusively on external programs and
        files, with no access to shell state:
        roughly 450 lines.
\end{itemize}

The third category is the refactor target. It represents
approximately 60\% of the file by line count and is the source of
the file's loading cost, its opacity to static analysis, and its
unsuitability for cross-host reuse.

\subsection{The script directory}

\texttt{\textasciitilde/bin} holds approximately sixty top-level
entries. Distribution by category:

\begin{itemize}
  \item Hardened scripts with shebangs, quoted variables, and
        meaningful exit codes (\texttt{aws\_*.sh},
        \texttt{rclone\_*.sh}, \texttt{git-*.sh}, and similar):
        roughly 25 entries.
  \item Microscripts under fifteen lines that reduce to a single
        command pipeline (the alias-in-disguise category):
        4 entries.
  \item Microscripts that are correctly placed but unhardened (no
        shebang, backticks, unquoted variables, dead code):
        5 entries.
  \item Non-executables (markdown documentation, PDFs, R package
        metadata, project subdirectories): approximately 12 entries
        plus 8 subdirectories.
  \item A buried \texttt{archive/} subdirectory containing 40+
        legacy files dating from 2021--2024, with extensive
        duplication.
\end{itemize}

\section{Guiding Principle}

The rule that drives the refactor is one sentence:

\begin{quote}
A helper should be a shell function only when it must mutate the
calling shell. Otherwise it should be a versioned, executable script
in \texttt{\textasciitilde/bin}.
\end{quote}

The four operations that require the function form, and that no
script can perform on the parent shell's behalf, are:

\begin{enumerate}
  \item Change the working directory.
  \item Export, set, or unset shell variables.
  \item Modify aliases, options, key bindings, or completion
        definitions.
  \item Manipulate the directory stack, the command history, or the
        job table.
\end{enumerate}

Anything outside that list (running a program, reading or writing
files, formatting output, fetching from a URL, calling
\texttt{git}, calling \texttt{make}, calling a database, parsing
JSON) communicates with the surrounding world through file
descriptors and exit codes alone, and is therefore a script's job.
The mechanical reason is that a script runs in a child process
spawned by \texttt{fork} and \texttt{execve}: the child receives a
copy of the parent's exported environment, file descriptors, and
arguments, but cannot modify the parent's memory. When the child
exits, the parent learns only the bytes written to \texttt{stdout},
the bytes written to \texttt{stderr}, and an integer exit code.

The single explicit exception is \texttt{source} (equivalently
\texttt{.}): the sourced file is read into the current shell rather
than run in a child, so a \texttt{cd} or \texttt{export} inside a
sourced file does affect the parent. Sourcing is reserved for
files whose authors explicitly intend to operate on the parent's
state (the startup file itself; the conda hook; plugin loaders);
ordinary helpers should be invoked, not sourced.

\section{Phased Refactor}

Each phase below is independently shippable. Execution may stop
between any two phases without leaving the toolbox in a worse
state than at the start.

\subsection{Phase 1: extract the largest function}

The dominant function by line count is \texttt{gz()}, an
interactive git-commit workflow with secret scanning, together
with its private helper \texttt{\_gz\_secret\_scan()} and a
companion \texttt{gz-scan-history()}. Combined size: roughly 330
lines spanning startup-file lines 267--597. The function calls
\texttt{git}, \texttt{gitleaks}, \texttt{fzf}, \texttt{awk}, and
\texttt{sed}; it reads from \texttt{stdin}; it writes to
\texttt{stdout}, \texttt{stderr}, and a log file. It mutates no
shell state.

\paragraph{Action.}
\begin{itemize}
  \item Create \texttt{\textasciitilde/bin/gz} with shebang
        \texttt{\#!/usr/bin/env zsh}. The function uses
        zsh-specific constructs (\texttt{\$\{(f)...\}},
        \texttt{\$\{arr:\#\}}, \texttt{\$\{(j:,:)...\}},
        \texttt{\$\{=var\}}, \texttt{typeset -U}) that bash will
        not parse.
  \item Inline \texttt{\_gz\_secret\_scan} into the same file (or
        place it in \texttt{\textasciitilde/bin/lib/} and source
        it).
  \item Create \texttt{\textasciitilde/bin/gz-scan-history} with
        the same shebang.
  \item Apply hardening: \texttt{set -uo pipefail}. (Omit
        \texttt{set -e} because both helpers conditionally probe
        for gitleaks subcommands and tolerate non-zero exits in
        that probe.)
  \item Resolve all \texttt{shellcheck} findings or annotate
        per-line with a justified \texttt{disable} comment.
  \item Delete the three function definitions from \texttt{.zshrc}.
  \item Verify by running \texttt{gz} end-to-end on a scratch
        repository, including the secret-scan refusal path and the
        \texttt{GZ\_ALLOW\_SECRET\_OVERRIDE} audited-override path.
\end{itemize}

\paragraph{Effort.} Half a day. Most of the time goes to the
end-to-end verification, since the function has multiple branches
that must each be exercised.

\paragraph{Acceptance criteria.}
\begin{itemize}
  \item \texttt{gz} and \texttt{gz-scan-history} run from a fresh
        shell without sourcing the dotfile.
  \item \texttt{shellcheck \textasciitilde/bin/gz} returns clean.
  \item Shell startup time, measured by
        \texttt{time zsh -i -c exit}, decreases by a measurable
        margin (typically 30--80\,ms on this size of extraction).
\end{itemize}

\subsection{Phase 2: extract the small functions}

After Phase~1, several smaller functions remain in \texttt{.zshrc}
that also fail the principle:

\begin{longtable}{@{}lp{3.0in}l@{}}
  \toprule
  \textbf{Function} & \textbf{Behaviour} & \textbf{Destination} \\
  \midrule
  \texttt{pp()}   & \texttt{rg --files | rg '\textbackslash.pdf\$' | fzf | xargs zathura}                    & \texttt{\textasciitilde/bin/pp}   \\
  \texttt{rr()}   & \texttt{rg --files | rg '\textbackslash.(R|Rmd)\$' | fzf | xargs vim}                    & \texttt{\textasciitilde/bin/rr}   \\
  \texttt{mma()}  & one-arg wrapper around the Mathematica kernel binary                                     & \texttt{\textasciitilde/bin/mma}  \\
  \texttt{mr()}   & \texttt{make -C "\$root" "\$\{@:-r\}"} from any project subdirectory                     & \texttt{\textasciitilde/bin/mr}   \\
  \texttt{nav()}  & prints a static help table of the navigation shortcuts                                   & \texttt{\textasciitilde/bin/nav}  \\
  \bottomrule
\end{longtable}

\paragraph{Action.} For each helper:

\begin{itemize}
  \item Copy the body to a file under
        \texttt{\textasciitilde/bin/}, mark executable.
  \item Add a shebang (\texttt{\#!/usr/bin/env bash} for the four
        bash-portable cases; \texttt{\#!/usr/bin/env zsh} only if
        the body uses zsh-only constructs).
  \item Add \texttt{set -euo pipefail} unless a deliberate
        deviation is documented in a comment.
  \item Quote variable expansions; replace any backticks with
        \texttt{\$(...)}.
  \item Resolve \texttt{shellcheck} findings.
  \item Delete the function definition from \texttt{.zshrc}.
\end{itemize}

\paragraph{Effort.} One hour total across all five helpers.

\paragraph{Acceptance criteria.} Each script is callable from a
fresh shell and from a non-interactive context (for example, from
within a Makefile). Each is \texttt{shellcheck}-clean.

\paragraph{Note on \texttt{mr}.} The \texttt{mr} function calls
\texttt{\_zzcollab\_root}, which lives in \texttt{.zshrc} as a
function and is not exported to child processes. The extracted
\texttt{mr} script must therefore inline a small project-root
walk rather than call the function. A few self-contained lines
at the top of the script suffice.

\subsection{Phase 3: promote microscripts to aliases}

Four scripts in \texttt{\textasciitilde/bin} reduce to a single
command pipeline with no flag parsing or branching. Each is more
honestly expressed as an alias.

\begin{longtable}{@{}lp{2.7in}l@{}}
  \toprule
  \textbf{Script} & \textbf{Body} & \textbf{Action} \\
  \midrule
  \texttt{\textasciitilde/bin/gg}    & \texttt{nohup .../ghostty >/dev/null 2>\&1 \&}             & alias and delete \\
  \texttt{\textasciitilde/bin/k}     & \texttt{nohup .../kitty >/dev/null 2>\&1 \&}              & alias and delete \\
  \texttt{\textasciitilde/bin/bb}    & \texttt{cd \textasciitilde/docs/personal; vim preludelist.tex} & alias as \texttt{vim} with absolute path; delete \\
  \texttt{\textasciitilde/bin/start} & \texttt{git-repo-push.sh; brew upgrade; brew doctor; brew cleanup} & alias OR keep as a hardened script; author's choice \\
  \bottomrule
\end{longtable}

\paragraph{Action.} Add the alias definitions to the relevant
section of \texttt{.zshrc}. Delete the corresponding scripts.
Confirm that no Makefile, launchd plist, or other repository
references the deleted file names.

\paragraph{Note on \texttt{bb}.} The current implementation
performs an unnecessary \texttt{cd} before invoking \texttt{vim}.
The \texttt{cd} happens in the script's child process and is
discarded on exit (it has no observable effect). The cleanest
alias passes the absolute path directly to \texttt{vim}:
\texttt{alias bb='vim \textasciitilde/docs/personal/preludelist.tex'}.

\paragraph{Effort.} Fifteen minutes.

\paragraph{Acceptance criteria.} \texttt{alias gg}, \texttt{alias
k}, \texttt{alias bb}, and (if chosen) \texttt{alias start} resolve
to the expected definitions. The four files are absent from
\texttt{\textasciitilde/bin}.

\subsection{Phase 4: harden the remaining microscripts}

Five scripts (\texttt{qcp}, \texttt{ss}, \texttt{vv}, \texttt{zz},
\texttt{p2}) are correctly placed but unhardened. All five share
the same defects:

\begin{enumerate}
  \item No shebang. The \texttt{file} utility reports them as
        `ASCII text', not executables. They run only because the
        parent shell falls back to \texttt{sh}.
  \item Backticks for command substitution rather than
        \texttt{\$(...)}.
  \item Unquoted variable expansions (\texttt{\$filename},
        \texttt{\$1}, \texttt{\$proj}). Filenames containing
        whitespace will fail silently or, worse, succeed with the
        wrong file.
  \item In four of the five (\texttt{qcp}, \texttt{ss},
        \texttt{vv}, \texttt{zz}), a variable named \texttt{proj}
        is computed from \texttt{\$1} or \texttt{\$(pwd)} and then
        never used except in debug \texttt{echo} calls. This is
        dead code.
\end{enumerate}

\paragraph{Action.} Apply the standard hardening pass to each
script: shebang, \texttt{set -euo pipefail}, replace backticks,
quote variables, delete dead code, resolve all
\texttt{shellcheck} findings.

\paragraph{Effort.} One hour. Phase~5 reduces this further by
collapsing the worst three.

\paragraph{Acceptance criteria.}
\texttt{file \textasciitilde/bin/\{qcp,ss,vv,zz,p2\}} reports each
as a shell script. \texttt{shellcheck} returns clean.

\subsection{Phase 5: consolidate the latest-of helpers}

Three of the scripts hardened in Phase~4 (\texttt{ss},
\texttt{vv}, \texttt{zz}) are the same script with different
terminal commands at the bottom: each finds the most recent file
matching a glob in the current directory, then opens it with a
chosen viewer (\texttt{open -a Skim}, \texttt{vim},
\texttt{zathura}). The duplication is gratuitous.

\paragraph{Action.} Replace the three scripts with one
parameterised helper plus three thin wrappers. The replacement
(verbatim) is given in Appendix~A.

\paragraph{Effort.} Twenty minutes.

\paragraph{Acceptance criteria.} Each of \texttt{ss}, \texttt{vv},
and \texttt{zz} continues to do exactly what it did before. The
shared logic appears once. Adding a fourth viewer (for example,
\texttt{Preview.app}) requires one new two-line wrapper.

\subsection{Phase 6: clean up non-executables}

The \texttt{\textasciitilde/bin} directory contains the following
non-executable artefacts that should be relocated:

\begin{itemize}
  \item Documentation: \texttt{CLAUDE.md},
        \texttt{README\_zzcollab.md},
        \texttt{date\_filtering.md},
        \texttt{refactor\_plan.md}. Move to a documentation
        directory or to the relevant project repository.
  \item R package metadata: \texttt{DESCRIPTION},
        \texttt{LICENSE}, the \texttt{inst/} subdirectory, the
        \texttt{tests/} subdirectory. Move to the package's own
        repository.
  \item PDFs: \texttt{README\_zzcollab\_2026-02-03\_0556.pdf},
        \texttt{refactor\_plan\_2026-04-25\_0930.pdf}. Move to a
        reference archive.
  \item Project subdirectories: \texttt{analysis/},
        \texttt{aws/}, \texttt{data/}, \texttt{init/},
        \texttt{launchd/}. These appear to be entire project
        workspaces that wandered into \texttt{\textasciitilde/bin}.
        Confirm and relocate each.
\end{itemize}

\paragraph{Action.} For each item, identify the natural home and
move it. Update any cross-references (Makefiles, launchd plists,
documentation). Run \texttt{ls \textasciitilde/bin} after the
sweep to confirm the directory now contains only executables and,
optionally, a single \texttt{lib/} subdirectory for sourced
helpers.

\paragraph{Effort.} One hour.

\paragraph{Acceptance criteria.} \texttt{find \textasciitilde/bin
-maxdepth 1 -not -type f -not -name 'lib' -not -name '.'} returns
nothing (or only the optional \texttt{lib/} directory).

\subsection{Phase 7: decommission the archive directory}

The buried \texttt{\textasciitilde/bin/archive/} subdirectory
contains 40+ legacy files dating from 2021--2024, including
multiple obvious duplicates (\texttt{vv}/\texttt{vvv}/\texttt{vvv.R};
\texttt{sm}/\texttt{sm.R}/\texttt{sm.rmd}/\texttt{sm.txt};
\texttt{tt}/\texttt{ttt}; \texttt{nn}/\texttt{nnn};
\texttt{rrtools\_original.sh}/\texttt{rrtools.sh.bak2}); broken
symlinks (\texttt{firefox.macos -> ff},
\texttt{mvim -> /usr/bin/gvimdiff}); and stub files of fewer than
fifty bytes.

The directory exists because at some prior point the author
hesitated to delete the files and stowed them out of sight
instead. That decision should be revisited now that the dotfiles
are under version control: git history preserves the removed
code, and the working tree should not.

\paragraph{Action.} Delete the directory in a single commit with
a descriptive message that names the rationale (`legacy
unreferenced helpers, last touched 2024-XX-XX, removed during
toolbox refactor; recoverable from git history if ever needed').

\paragraph{Effort.} Five minutes (plus a brief audit to confirm
none of the archived files is referenced by a script or
scheduler still in use).

\paragraph{Acceptance criteria.} \texttt{\textasciitilde/bin/archive/}
no longer exists. \texttt{git log} preserves the deleted
content for any future recovery.

\subsection{Phase 8: naming hygiene}

Several helpers carry names that fail the legibility standard the
rest of the refactor enforces. The renames proposed below are for
review, not for unilateral application: short names are sometimes
defensible on grounds of frequency-of-use, and the author is the
only person who knows which fall into that category. The proposed
disposition is divided into three sub-phases by category.

\subsubsection*{8a. Application launchers}

Two scripts launch terminal emulators with single- or
double-character names that collide with conventional grep
targets (\texttt{gg} is also common shorthand for
\texttt{git grep}). After Phase~3 these become aliases; the
rename and the promotion should land in the same commit.

\begin{longtable}{@{}llp{2.7in}@{}}
  \toprule
  \textbf{Current} & \textbf{Proposed} & \textbf{Rationale} \\
  \midrule
  \texttt{\textasciitilde/bin/gg} & \texttt{alias ghostty=...} & Matches the binary name \\
  \texttt{\textasciitilde/bin/k}  & \texttt{alias kitty=...}   & Matches the binary name \\
  \bottomrule
\end{longtable}

\subsubsection*{8b. Gratuitously cryptic helpers}

The following entries in \texttt{\textasciitilde/bin} have names
that do not map to their behaviour, are infrequently invoked,
and (in several cases) cannot be readily recalled without
opening the file. Each requires a single-sitting audit ending in
either a descriptive rename plus a one-line header comment, or
removal.

\begin{longtable}{@{}llp{2.5in}@{}}
  \toprule
  \textbf{Helper} & \textbf{Behaviour} & \textbf{Disposition} \\
  \midrule
  \texttt{bb}    & edits \texttt{\textasciitilde/docs/personal/preludelist.tex} & rename to \texttt{edit-preludes}, or absorb as the Phase 3 alias \\
  \texttt{qcp}   & copies a built post PDF into the working directory          & rename to \texttt{pull-post-pdf}                                  \\
  \texttt{p2}    & enscript-to-PDF-to-printer pipeline (two-up)                 & rename to \texttt{print-2up}                                      \\
  \texttt{ss}    & opens latest PDF in Skim                                     & retire after Phase~5 (replaced by \texttt{latest-of} wrapper)     \\
  \texttt{vv}    & edits latest Rmd/qmd in vim                                  & retire after Phase~5                                              \\
  \texttt{zz}    & opens latest PDF in zathura                                  & retire after Phase~5                                              \\
  \texttt{n1}    & unknown (33 lines)                                           & inspect; rename to a descriptive name or delete                   \\
  \texttt{n2}    & unknown (22 lines)                                           & inspect; rename or delete                                         \\
  \texttt{p1}    & unknown (23 lines)                                           & inspect; rename or delete                                         \\
  \texttt{tn}    & unknown (546 lines, the largest unnamed helper)              & inspect; rename or delete; if kept, add a header comment         \\
  \texttt{wisu}  & undocumented Mach-O binary                                   & inspect; if author cannot recall purpose, delete and recompile from source if ever needed again \\
  \bottomrule
\end{longtable}

\paragraph{Action.} For each rename:
\begin{enumerate}
  \item Search for cross-references before touching the file:
        \texttt{grep -r 'old-name' \textasciitilde/Library/LaunchAgents
        /etc/cron* \$DOTFILES} (adjust paths to environment).
  \item Use \texttt{git mv old new} so that history is preserved.
  \item Update any references found in the previous step in the
        same commit as the rename.
  \item Add a one-line header comment to any helper kept under
        an unfamiliar name, summarising what it does.
\end{enumerate}

\paragraph{Effort.} One hour, dominated by the audit of the
unknown helpers. The renames themselves are individually cheap.

\paragraph{Acceptance criteria.}
\begin{itemize}
  \item Every entry in \texttt{\textasciitilde/bin} either has a
        name a reader can guess from its purpose, or has a
        one-line header comment that clarifies the purpose.
  \item No file in \texttt{\textasciitilde/bin} produces `I do
        not remember what this does' on inspection.
  \item All scheduler entries (launchd, systemd, cron, Make)
        reference the new names.
\end{itemize}

\subsection{Phase 9: install a guardrail}

The refactor of phases 1--8 cleans up a decade of accreted drift.
Without an automated check, the same drift will reaccumulate
within a few years.

\paragraph{Action.} Add a \texttt{Makefile} target to the
dotfiles repository:

\begin{lstlisting}[language=make]
lint:
	@find ~/bin -maxdepth 1 -type f -perm +111 \
	    | xargs -I{} sh -c \
	      'head -1 "$$1" | grep -q "^#!.*sh" && echo "$$1"' _ {} \
	    | xargs shellcheck
\end{lstlisting}

Wire \texttt{make lint} into a pre-commit hook on the dotfiles
repository so that any new helper added to
\texttt{\textasciitilde/bin} is statically analysed before it
lands in version control. The check is intentionally narrow: it
covers only shell scripts and only the top-level
\texttt{\textasciitilde/bin}, not the long tail of subdirectories
that may hold helpers in other languages.

\paragraph{Effort.} Thirty minutes (including hook installation
and a verification pass against the post-Phase-6 toolbox).

\paragraph{Acceptance criteria.} \texttt{make lint} returns clean
on the post-refactor toolbox. A deliberately broken script
introduced as a test causes the pre-commit hook to fail.

\section{Effort and Sequencing}

\begin{longtable}{@{}lll@{}}
  \toprule
  \textbf{Phase} & \textbf{Description}                                  & \textbf{Effort} \\
  \midrule
  1 & Extract \texttt{gz}, \texttt{\_gz\_secret\_scan}, \texttt{gz-scan-history} & 4\,h            \\
  2 & Extract \texttt{pp}, \texttt{rr}, \texttt{mma}, \texttt{mr}, \texttt{nav}  & 1\,h            \\
  3 & Promote four microscripts to aliases                                       & 0.3\,h          \\
  4 & Harden five remaining microscripts                                         & 1\,h            \\
  5 & Consolidate \texttt{ss}, \texttt{vv}, \texttt{zz} into one helper          & 0.3\,h          \\
  6 & Relocate non-executables out of \texttt{\textasciitilde/bin}               & 1\,h            \\
  7 & Decommission \texttt{\textasciitilde/bin/archive/}                         & 0.1\,h          \\
  8 & Naming hygiene (launchers, cryptic helpers)                                & 0.7\,h          \\
  9 & Install \texttt{make lint} guardrail                                       & 0.5\,h          \\
  \midrule
    & \textbf{Total}                                                             & \textbf{8.9\,h} \\
  \bottomrule
\end{longtable}

The phases are ordered by leverage: the early phases produce most
of the visible improvement. A single working day suffices for the
full plan, although the natural cadence is a phase per evening
across two weeks, with each commit landed independently and
verified in normal use before the next phase begins.

\section{Risks and Mitigations}

\paragraph{Schedulers refer to file names.}
If a launchd plist, systemd unit, cron entry, or Makefile
references a script by name, renaming or deleting the script
silently breaks the scheduler. Before any rename or delete, run
\texttt{grep -r "script-name" \textasciitilde/Library/LaunchAgents
\textasciitilde/etc/cron* \$DOTFILES} (adjust paths to your
environment). Land both the rename and the scheduler update in
the same commit.

\paragraph{Zsh-only constructs in extracted code.}
\texttt{gz} and possibly other extractions use zsh-specific
syntax that bash will not parse. The shebang must match the
syntactic features actually in use. A bash shebang on a function
that uses \texttt{\$\{(f)...\}} fails in unhelpful ways; default
to \texttt{\#!/usr/bin/env zsh} for any extraction whose syntax
has not been explicitly portabilised.

\paragraph{Working directory and PATH inheritance.}
A script extracted from a function inherits \texttt{PATH} and
\texttt{\$PWD} from the calling shell, which is generally
correct. The exception arises when the parent function previously
ran inside a block that prepended a directory to \texttt{PATH};
the new script needs to do its own \texttt{PATH} adjustment or
use absolute paths. Verify by running the extracted script with
\texttt{env -i} to surface any implicit-environment dependency.

\paragraph{Secrets in extracted code.}
A function previously sourced from \texttt{.zshrc} may have
benefited from environment variables set elsewhere in the same
file (the pattern \texttt{[[ -f \textasciitilde/.env ]] \&\& source
\textasciitilde/.env} appears at line 25 of the audited file).
The extracted script does not inherit those variables unless they
were already \texttt{export}ed. Either source the env file
explicitly inside the script or read secrets via a credential
helper such as \texttt{pass}.

\paragraph{Shellcheck false positives.}
External programs invoked with computed argument arrays
(\texttt{"\$\{args[@]\}"}) sometimes trigger \texttt{SC2068},
\texttt{SC2086}, or \texttt{SC2154}. Disable per-line with a
comment that names the reason; never blanket-disable a warning
across the project.

\paragraph{Loss of in-session state on accidental sourcing.}
A sourced file shares the parent's namespace and can clobber any
name defined there. An error in a sourced file (an unset variable
under \texttt{set -u}, a stray \texttt{exit}) terminates the
interactive session rather than a child process. The phase~1--6
extractions are designed to be invoked, not sourced; verify after
each extraction that no stale documentation or muscle memory is
\texttt{source}ing them.

\section{Validation Plan}

After each phase, verify the following before proceeding to the
next:

\begin{enumerate}
  \item \textbf{Behaviour.} The affected helpers continue to
        produce identical observable output, exit codes, and
        side effects on a representative input set.
  \item \textbf{Lint.} \texttt{shellcheck \$f} returns clean for
        each touched file in \texttt{\textasciitilde/bin}.
  \item \textbf{Startup time.} \texttt{time zsh -i -c exit}
        executed five times produces a stable startup time at or
        below the post-Phase-1 baseline.
  \item \textbf{Discoverability.} Each new helper in
        \texttt{\textasciitilde/bin} appears under tab completion
        against \texttt{\$PATH}; each new alias appears in
        \texttt{alias} output.
  \item \textbf{Reversibility.} The phase's commit is in version
        control on its own. Reverting the single commit restores
        the prior behaviour without manual cleanup.
\end{enumerate}

A short script (\texttt{\textasciitilde/bin/dotfiles-doctor})
that runs all five checks and reports a single pass/fail line
is a useful by-product of the refactor.

\section{Expected Outcomes}

After execution of all nine phases:

\begin{itemize}
  \item \texttt{.zshrc} shrinks from approximately 736 lines to
        approximately 350 lines. The remainder is OS-conditional
        configuration, plugin and completion initialisation, the
        prompt, and the legitimate shell-state functions.
  \item \texttt{\textasciitilde/bin} becomes a flat directory of
        correctly-shebanged, \texttt{shellcheck}-clean
        executables (with at most a \texttt{lib/} subdirectory
        for sourced helpers).
  \item Shell startup time decreases by an estimated 50--150\,ms
        from the function extractions in phases 1--2. The
        dominant contribution comes from removing the \texttt{gz}
        family.
  \item Every helper in the toolbox is one \texttt{git log} away
        from full history, one \texttt{shellcheck} away from
        lint, and one \texttt{cat} away from a complete reading.
  \item The \texttt{make lint} guardrail prevents the same
        drift from reaccumulating.
\end{itemize}

\paragraph{Caveat on the startup-time projection.} The 50--150\,ms
estimate above accounts only for the function-extraction phases.
After execution, the dominant remaining startup cost in the
audited file is the conda initialisation block (lines 634--647),
which invokes \texttt{conda 'shell.zsh' 'hook'} on every shell
launch, spawning a Python interpreter. On the audited machine
this adds roughly 50--150\,ms by itself, so the perceived
improvement to interactive responsiveness will be smaller than
the function-extraction numbers alone suggest. The standard
mitigation is to lazy-load conda: replace the eager hook with a
shell stub that initialises conda only on first \texttt{conda}
invocation. This is out of scope for the current plan but worth
considering as a follow-up; it is the highest-leverage remaining
intervention on shell startup time.

The largest qualitative gain is composability. Scripts in
\texttt{\textasciitilde/bin} are picked up by anything that
respects \texttt{PATH}: cron, launchd, Make rules, Docker
containers that mount the home directory, HPC submission
scripts, scheduled GitHub Actions runners. Functions buried in
\texttt{.zshrc} are visible only to interactive zsh sessions on
the host that holds the file. The phases above convert helpers
that the author has already written into helpers that any
tooling can call.

\appendix

\section{Replacement for ss / vv / zz (Phase 5)}

Place at \texttt{\textasciitilde/bin/latest-of}:

\begin{lstlisting}[language=bash]
#!/usr/bin/env bash
# latest-of: open the most recent file in $PWD whose name matches
# a regex, using the named viewer. Designed to be wrapped by
# small helpers (see ss, vv, zz).
#
# Usage: latest-of <regex> <viewer> [viewer-args...]
set -euo pipefail

pattern="${1:?usage: latest-of <regex> <viewer> [args...]}"
viewer="${2:?usage: latest-of <regex> <viewer> [args...]}"
shift 2

# ls -At sorts by mtime, newest first; head -n 1 picks the newest
# match. grep -E supplies extended regex.
file=$(ls -At | grep -E -- "$pattern" | head -n 1 || true)
if [[ -z "$file" ]]; then
    echo "latest-of: no file in $PWD matches /$pattern/" >&2
    exit 1
fi

exec "$viewer" "$@" -- "$file"
\end{lstlisting}

Then replace \texttt{ss}, \texttt{vv}, \texttt{zz} with two-line
wrappers (or aliases):

\begin{lstlisting}[language=bash]
#!/usr/bin/env bash
# ~/bin/ss: open the most recent PDF with Skim.
exec latest-of '\.pdf$' open -a Skim
\end{lstlisting}

\begin{lstlisting}[language=bash]
#!/usr/bin/env bash
# ~/bin/vv: edit the most recent R/Markdown document with vim.
exec latest-of '\.(R|r)?md$|\.qmd$' vim
\end{lstlisting}

\begin{lstlisting}[language=bash]
#!/usr/bin/env bash
# ~/bin/zz: open the most recent PDF with zathura.
exec latest-of '\.pdf$' zathura
\end{lstlisting}

\section{Glossary}

\begin{itemize}[leftmargin=1em,itemsep=4pt]
  \item \textbf{Process.}
    An operating-system execution unit with its own memory,
    environment, file descriptors, and a unique process
    identifier. The operating system isolates processes so that
    one cannot read or write another's memory directly.
  \item \textbf{State.}
    The information a process holds in memory at a given moment.
    For a shell process, state includes the working directory,
    the values of environment and local variables, defined
    aliases and functions, shell options, the history list, the
    directory stack, and the job table.
  \item \textbf{Namespace.}
    The set of names (variables, functions, aliases) visible to
    a particular execution context. The interactive shell
    maintains one namespace; a child process receives a copy of
    the exported portion and nothing else.
  \item \textbf{Scope.}
    The region of code in which a given name is defined and
    accessible. Functions and scripts have distinct scopes by
    construction.
  \item \textbf{Reproducibility.}
    The property that re-running the same code on the same
    inputs in the same environment yields the same outputs.
    Requires both that the code be captured and that the
    environment be characterised.
  \item \textbf{Source (verb).}
    To execute a file's commands in the current shell rather
    than in a child process; the only mechanism by which a file
    can mutate the calling shell's state.
\end{itemize}

\vfill
\hrule
\vspace{0.5em}
{\small \emph{Rendered on 2026-05-03 at 11:14 PDT.}\\
\emph{Source:}
\texttt{\textasciitilde/prj/qblog/posts/51-scripts-vs-functions/}\\
\hspace*{1.5em}\texttt{scripts-vs-functions/analysis/configs/}\\
\hspace*{1.5em}\texttt{dotfiles\_refactor\_plan.tex}}

\end{document}
